\begin{explanationbox}
	\textbf{\textcolor{CExplanation}{一句话直觉}}

	正三棱柱的外接球球心位于上下底面中心连线的中点，半径由底面正三角形的外接圆半径和高的一半通过勾股定理决定。

	\medskip
	\textbf{\textcolor{CExplanation}{核心拆解}}
	\begin{description}[style=nextline, leftmargin=2.8em, labelsep=0.5em, itemsep=0.45em]
		\item[\textbf{要点一（球心位置）：}] 球心在上下底面中心连线上，并且位于这条线段的中点.
		\item[\textbf{要点二（底面半径）：}] 底面正三角形边长为 $a$，其外接圆半径为 $r=a/\sqrt{3}$.
		\item[\textbf{要点三（半径构成）：}] 外接球半径 $R$ 满足
		      \[
			      R^{2}=r^{2}+\left(\frac h2\right)^{2}.
		      \]
	\end{description}

	\medskip
	\textbf{\textcolor{CExplanation}{几何本质（对称与勾股）}}

	正三棱柱上下两个底面全等且平行，对应顶点关于中截面对称。外接球球心到所有顶点距离相等，因此球心位于上下底面中心连线的中点。连接球心、上底面中心和上底面一个顶点，可以得到一个直角三角形，于是
	\[
		R^{2}=r^{2}+\left(\frac h2\right)^{2}.
	\]

	\medskip
	\textbf{\textcolor{CExplanation}{代数意义（分离变量）}}

	公式将底面条件（$a$）与高度条件（$h$）分离在 $a^{2}/3$ 和 $h^{2}/4$ 两项中，不产生交叉项，便于分别处理。

	\medskip
	\textbf{\textcolor{CExplanation}{考点价值（公式应用）}}
	\begin{itemize}[leftmargin=*, itemsep=0.5em, topsep=0.4em]
		\item \textbf{考法一（直接计算）：} 已知 $a$ 和 $h$，直接代公式求 $R$，进而求外接球的表面积、体积.
		\item \textbf{考法二（反求参数）：} 已知 $R$ 与其中一个量，通过方程解出另一个量.
		\item \textbf{考法三（判断存在性）：} 反求参数时还要检查正数条件：若反求 $a$，需
		      \[
			      R^{2}>\frac{h^{2}}{4};
		      \]
		      若反求 $h$，需
		      \[
			      R^{2}>\frac{a^{2}}{3}.
		      \]
		      否则不存在满足条件的正三棱柱.
	\end{itemize}

	\medskip
	\textbf{\textcolor{CExplanation}{推广视角（只限直棱柱）}}

	对于一般直三棱柱，若底面三角形的外接圆半径为 $R_{\triangle}$，棱柱高为 $h$，则有
	\[
		R^{2}=R_{\triangle}^{2}+\left(\frac h2\right)^{2}.
	\]
	正三棱柱只是其中的特殊情况，因为底面正三角形的外接圆半径恰好为 $a/\sqrt3$。但斜棱柱不能直接套用本结论。

	\medskip
	\textbf{\textcolor{CExplanation}{顿悟点}}

	\textbf{上下底面中心连线的中点是外接球球心；底面正三角形的外接圆半径是 $a/\sqrt3$，这两个点抓住，公式就自然出来了。}

	\medskip
	\textbf{\textcolor{CExplanation}{使用场景}}
	\begin{itemize}[leftmargin=*, itemsep=0.5em, topsep=0.4em]
		\item \textbf{场景一（直接应用）：} 给出正三棱柱的 $a$ 和 $h$，直接利用公式求 $R$.
		\item \textbf{场景二（建立方程）：} 已知外接球半径 $R$ 和部分条件，通过公式建立方程求解.
		\item \textbf{场景三（防止误用）：} 遇到底面不是正三角形或侧棱不垂直底面的三棱柱时，不能把 $a/\sqrt3$ 和 $h/2$ 直接硬套.
	\end{itemize}
\end{explanationbox}